\section{Related Work}
\subsection{Decomposition of Observed Travel Times}
Sparse probe observations can span several road links, requiring the elapsed time between observations to be allocated along an identified path. Hellinga et al.\ separate free-flow, congestion, and stopping time using link geometry and free-flow speeds\cite{hellinga2008}. Their congestion likelihood uses the preceding polling interval, while the stopping likelihood reflects the vehicle's position relative to the downstream end of a link. The resulting allocation integrates over possible congestion levels. Hofleitner and Bayen instead derive link-time distributions from a statistical model of arterial queues and driving behavior\cite{hofleitner2011}. These distributions are mixtures of log-concave components. The mixture allocation problem is nonconvex, but fixed delay patterns yield convex subproblems that support enumeration and EM-based algorithms. Both studies constrain allocated times to the observed duration.

Network-level inference also uses observations shared across trips. Yin et al.\ infer link-time distributions from entry and exit records, using Gaussian likelihoods, mixture estimation when paths are unknown, and trip splitting for more general distributions\cite{yin2015entryexit}. PR-LTTE alternates path recovery from trip distance and duration with nonnegative least-squares link-time estimation\cite{sun2022prltte}. These methods estimate shared road-link quantities and may recover missing paths. The present study assumes an ordered sequence of directed bus inter-stop intervals, each of which may span several road links. Its shared estimates vary with observed context, and its redistribution varies with the trip. The NNLS-Ridge reference examines fixed interval-identity coefficients under this known-sequence setting.

\subsection{Learned Travel-Time Representations}
DeepTTE combines geographic convolution and recurrent layers with joint whole-path and local-path supervision\cite{wang2018deeptte}. Its local targets are timestamp differences over GPS windows, which can overlap; the whole-path prediction is produced by a separate attention-based head. DutyTTE addresses uncertainty when the route is not supplied, combining predicted-path alignment with contextual segment uncertainty\cite{mao2025dutytte}. These designs show how local observations and route representations can support prediction of an unknown trip duration.

ProbETA provides a closer connection to sharing information across repeated links\cite{xu2025probeta}. Xu et al.\ model link travel time as a shared mean plus day-specific and trip-specific Gaussian effects. Learned low-rank covariance representations express dependencies both across trips and within a trip. Training uses joint likelihoods of trip durations and timestamp-derived sub-trip durations; inference conditions the queried trip's distribution on nearby completed trips. This is a probabilistic use of shared link evidence. In our setting, the queried total is already supplied, and limited singleton labels supervise the allocation. The trip correction is bounded relative to a shared contextual base and sums to zero within that trip, rather than parameterizing a covariance distribution.

\subsection{Aggregate Supervision and Consistent Outputs}
Learning from aggregate observations formalizes the estimation of individual labels from group-level supervision. Zhang et al.\ derive a likelihood for the observed aggregate and study regression from means or sums under conditional independence assumptions\cite{zhang2020aggregate}. Their aggregate error trains the individual predictor. Here, proportional rescaling enforces the supplied total for every output, so a loss on that final sum provides no learning signal. Training and model selection therefore use the visible singleton labels in the training and validation sets, respectively.

Forecast reconciliation offers another way to enforce additive consistency. MinT combines forecasts at different aggregation levels using their error covariance to obtain coherent forecasts\cite{wickramasuriya2019mint}. The present total is an observed input, while the unknown interval shares are learned and then scaled to it. We accordingly compare the shared-base redistribution model with direct interval prediction under the same inputs, label budgets, and final scaling. This comparison examines how the within-trip allocation is learned once consistency is imposed.
